\chapter{绪论}

\section{课题背景和意义}
本科毕业设计既是对专业知识的综合运用，也是对研究过程和工程表达能力的集中检验。一个可复现的论文模板能够把学校的格式要求落实为源文件中的稳定规则，减少手工调整带来的版式漂移。

\subsection{研究背景}
传统文字处理模板通常依赖本地软件状态。字体、目录域、交叉引用和页码更新如果没有被持续检查，提交前容易出现不一致。LaTeX 将结构和排版分开，使章节、公式、图表和参考文献可以由源文件重建\cite{lamport1994}。

\subsection{研究意义}
本项目提供一个面向江苏大学理工农医教类毕业设计（论文）的起点工程。使用者只需要修改元数据和章节文件，就可以生成与学校模板约束相符的初稿，并把注意力放在研究内容、数据和结论上。

\section{国内外相关工作}
学术论文模板通常包含文档类、字体配置、前置内容、章节组织、参考文献样式和构建脚本。相关工作表明，模板应同时提供清晰的默认值和可覆盖的接口，避免使用者直接修改类文件\cite{mittelbach2004}。本文以模块化源文件和自动化构建为核心，保留后续加入学科专用命令的空间。

\section{本文主要工作}
本文的主要工作包括：提取学校 2026 版通用格式中的可执行约束；实现 `ujsthesis` 文档类；组织封面、声明、摘要、目录、正文、参考文献、附录和致谢；最后通过 XeLaTeX、Biber 和 PDF 文本检查验证样例工程。
